% =============================================================================
%  Conclusion — synthèse de la campagne, couverture, suites
% =============================================================================
\section*{Conclusion}
\addcontentsline{toc}{section}{Conclusion}

\paragraph{État de la tour} Ce rapport présente la formalisation \emph{selon la
tour des théories} : la théorie logique (Partie~I) sert de base, la théorie des
ensembles (Partie~II) l'augmente de ses axiomes propres, et l'ordre, les
cardinaux puis les structures (Parties~III--IV) en sont des extensions. Une
trentaine de cibles ont été closes à travers les chapitres~II et~III
(correspondances, familles, produits, équivalences ; ordre, bon ordre,
cardinaux). Chaque résultat a franchi le même contrôle : preuve construite uniquement par les primitives
\Nstar{}, conclusion strictement égale (au renommage des liants $\tau$ près) à
l'énoncé de Bourbaki, hypothèses réduites aux antécédents honnêtes load-bearing,
invariant \texttt{theorie\_ensembles() == 22} préservé, et test miroir au vert.
La suite complète collecte sans erreur (plus de $3000$ tests).

\paragraph{Dossiers-trous comblés} La campagne a rempli plusieurs dossiers de
l'arborescence qui n'étaient que des marqueurs de manque : la diagonale
injective et son appartenance au produit $E^I$ (II.5.2/II.5.3), les ensembles
filtrants (III.1.8, élément maximal $=$ plus grand élément), et les parties
saturées (II.6.4). L'arborescence calquée sur le livre rend ces progrès
\emph{visibles} : un dossier jadis vide porte désormais un résultat certifié.

\paragraph{Honnêteté des statuts} Plusieurs résultats sont \emph{clos sous
hypothèses honnêtes} plutôt qu'inconditionnellement clos --- la réciproque de la
monotonie du produit (II.5.4) repose sur la surjectivité de la projection
$\mathrm{pr}_\alpha$, reçue en hypothèse ; l'associativité de la réunion (II.4.2)
sur les deux clauses fidèles modélisant $L = \bigcup_\lambda J_\lambda$. Dans
tous les cas la conclusion reste hors des hypothèses : ce sont des faits
conditionnels au sens de Bourbaki, non des troncatures. Les formes plus fortes
non atteintes (caractérisation par treillis de l'intersection d'intervalles,
forme cardinale complète) sont signalées comme résidus honnêtes.

\paragraph{Le mur cardinal} Les résultats d'arithmétique cardinale profonde
restent partiels, bloqués par un mur algorithmique (hachage de formules sur des
assemblages $\tau$ massifs) et non par une lacune mathématique. Le théorème de
Cantor $2^a > a$ a néanmoins été clos au niveau cardinal (III.3, Théorème~2) en
contournant la construction explicite par les lemmes d'équipotence déjà
certifiés.

\paragraph{Garde-fous méthodologiques} La campagne a été menée par délégation à
des agents, chaque preuve étant \emph{revérifiée par le noyau} avant
intégration : un agent ne peut pas forger un \theoreme{}. Les outils mutant le
dépôt sont transactionnels (préflight $+$ rollback) ; les commits sont
incrémentaux, sous garde anti-contamination ; toute exécution de test est bornée
par un \emph{timeout}.

\paragraph{Suites} Le moissonnage de l'audit de couverture identifie d'autres
cibles faisables (compléments du Théorème~1 sur les rétractions/sections, lois
de De~Morgan pour les familles, décomposition canonique d'une application,
propriétés order-théoriques du bon ordre) qui prolongeront le corpus selon le
même protocole. L'horizon reste la couverture intégrale du livre, les résultats
cardinaux profonds attendant un progrès sur le c\oe ur de substitution.

\paragraph{Vers une IA bâtisseuse de théories} Au-delà du livre certifié, la
finalité du projet est un \emph{substrat d'apprentissage} : le corpus consigne
non seulement les preuves, mais la \emph{construction} de chaque couche de la
tour, le \emph{pourquoi} de chaque choix, et les \emph{erreurs} corrigées
(annexe « Journal des erreurs \& leçons »). Une spécification fausse refusée par
le noyau, une hypothèse inerte retirée, une tautologie écartée, une capture de
variable déjouée : autant de données de première classe pour une intelligence
appelée, à terme, à construire des théories elle-même à partir d'axiomes ---
dans le respect de la même discipline de stratification et de certification.
